%--------------------------------------------------------------------------------
%
%
%--------------------------------------------------------------------------------
\chapter*{Synthetic Instructions}
\addcontentsline{toc}{chapter}{Synthetic Instructions}

The Twin-64 instruction set provides a rich set of options for individual instructions. Setting a defined option bit in an instruction can enable additional functionality with minimal impact on the overall data path. Similarly, instructions that use general register zero as an argument can take advantage of predefined behaviors. 

For better readability and convenience, an assembler can offer \emph{synthetic instructions}, which are higher-level mnemonics generated from the base instruction set. These synthetic instructions simplify programming by providing a more natural or concise representation of commonly used instruction patterns.

Synthetic instructions are also useful when a single instruction cannot fully express the desired operation. For example, loading an immediate value larger than the instruction field allows requires an instruction sequence. An assembler can provide a generic "load immediate" synthetic instruction that automatically emits such an instruction sequence depending on the size of the value. Similarly, if the offset to a base register is too small to reach a desired data item, the assembler can transparently emit an additional `ADDIL` instruction to compute the correct address.

\section*{Immediate Operations}
\addcontentsline{toc}{section}{Immediate Operations}

\begin{longtable}{@{}p{0.3\linewidth}p{0.3\linewidth}p{0.3\linewidth}@{}}
    % \caption{Immediate synthetic instructions} \\
    \toprule
    \textbf{OpCode} & \textbf{Sequence} & \textbf{Comment} \\
    \midrule
    \endfirsthead
    \toprule
    \textbf{OpCode} & \textbf{Sequence} & \textbf{Comment} \\
    \midrule
    \endhead
    \midrule
    \multicolumn{2}{r}{\textit{Continued on next page}} \\
    \midrule
    \endfoot
    \bottomrule
    \endlastfoot
    \textbf{LDI} Rn, val & LDO Rn, val(R0) & load an immediate value in the range of $\pm 2^{12}$ \\
	\midrule
	\textbf{LDI} Rn, val & LDIL Rn, L\%val \newline ADDIL Rn, R\%val & load an immediate value in the range of $\pm 2^{31}$ \\
	\midrule
	\textbf{LDI} Rn, val & LDIL Rn, L\%val \newline ADDIL Rn, R\%val \newline LDIL.M Rn, M\%val & load an immediate value in the range of $\pm 2^{51}$ \\
	\midrule
	\textbf{LDI} Rn, val & LDIL Rn, L\%val \newline ADDIL Rn, R\%val \newline LDI.M Rn, M\%val  \newline LDIL.U Rn, U\%val & load an immediate value in the range of $\pm 2^{63}$ \\
	\midrule
\end{longtable}

\section*{Register Operations}
\addcontentsline{toc}{section}{Register Operations}

\begin{longtable}{@{}p{0.3\linewidth}p{0.3\linewidth}p{0.3\linewidth}@{}}
    % \caption{Immediate synthetic instructions} \\
    \toprule
    \textbf{OpCode} & \textbf{Sequence} & \textbf{Comment} \\
    \midrule
    \endfirsthead
    \toprule
    \textbf{OpCode} & \textbf{Sequence} & \textbf{Comment} \\
    \midrule
    \endhead
    \midrule
    \multicolumn{2}{r}{\textit{Continued on next page}} \\
    \midrule
    \endfoot
    \bottomrule
    \endlastfoot
    \textbf{CLR} Rn & AND Rn, Rn, 0 & clears RegN. \\
    \midrule
    \textbf{CPY} Rn, Rm & OR Rn, Rm, 0 & copies RegM to RegN. \\
    \midrule
    \textbf{COM} Rn & XOR Rn, Rn, 1 & logically inverts RegN. \\
    \midrule
    \textbf{NEG} Rn & SUB Rn, R0, Rn & negates a value in RegN. \\
    \midrule
\end{longtable}

\section*{Arithmetic Operations}
\addcontentsline{toc}{section}{Arithmetic Operations}

\begin{longtable}{@{}p{0.3\linewidth}p{0.3\linewidth}p{0.3\linewidth}@{}}
    \toprule
    \textbf{OpCode} & \textbf{Sequence} & \textbf{Comment} \\
    \midrule
    \endfirsthead

    \toprule
    \textbf{OpCode} & \textbf{Sequence} & \textbf{Comment} \\
    \midrule
    \endhead

    \midrule
    \multicolumn{2}{r}{\textit{Continued on next page}} \\
    \midrule
    \endfoot

    \bottomrule
    \endlastfoot

    \textbf{INC} Rn, val  & ADD Rn, Rn, val & increments Rn by \textless val\textgreater . \\
    \midrule
    \textbf{DEC} Rn, val  & SUB Rn, Rn, val & decrements Rn by \textless val\textgreater . \\
    \midrule

    % Combine the next 3 rows:
    \textbf{EXT8}  Rn, Rm & EXTR.S Rn, Rm, 0, 8    & \multirow{3}{*}{sign extends a byte in Rn.} \\
    \textbf{EXT16} Rn, Rm & EXTR.S Rn, Rm, 0, 16   &  \\
    \textbf{EXT32} Rn, Rm & EXTR.S Rn, Rm, 0, 32   &  \\
    \midrule

\end{longtable}



\section*{Shift and Rotate Operations}
\addcontentsline{toc}{section}{Shift and Rotate Operations}

\begin{longtable}{@{}p{0.3\linewidth}p{0.3\linewidth}p{0.3\linewidth}@{}}
    % \caption{Immediate synthetic instructions} \\
    \toprule
    \textbf{OpCode} & \textbf{Sequence} & \textbf{Comment} \\
    \midrule
    \endfirsthead
    \toprule
    \textbf{OpCode} & \textbf{Sequence} & \textbf{Comment} \\
    \midrule
    \endhead
    \midrule
    \multicolumn{2}{r}{\textit{Continued on next page}} \\
    \midrule
    \endfoot
    \bottomrule
    \endlastfoot
    \textbf{ASR} & ... & ... \\
    \midrule
    \textbf{LSR} & ... & ... \\
    \midrule
    \textbf{ROR} & ... & ... \\
    \midrule
    \textbf{ROL} & ... & ... \\
    \midrule

\end{longtable}

\section*{Memory Reference Operations}
\addcontentsline{toc}{section}{Memory Reference Operations}

\begin{longtable}{@{}p{0.3\linewidth}p{0.3\linewidth}p{0.3\linewidth}@{}}
    % \caption{Immediate synthetic instructions} \\
    \toprule
    \textbf{OpCode} & \textbf{Sequence} & \textbf{Comment} \\
    \midrule
    \endfirsthead
    \toprule
    \textbf{OpCode} & \textbf{Sequence} & \textbf{Comment} \\
    \midrule
    \endhead
    \midrule
    \multicolumn{2}{r}{\textit{Continued on next page}} \\
    \midrule
    \endfoot
    \bottomrule
    \endlastfoot
    \textbf{LOAD} & ... & ... \\
    \midrule
    \textbf{STOR} & ... & ... \\
    \midrule
\end{longtable}

\section*{Branch Type Operations}
\addcontentsline{toc}{section}{Branch Type Operations}

\begin{longtable}{@{}p{0.3\linewidth}p{0.3\linewidth}p{0.3\linewidth}@{}}
    % \caption{Immediate synthetic instructions} \\
    \toprule
    \textbf{OpCode} & \textbf{Sequence} & \textbf{Comment} \\
    \midrule
    \endfirsthead
    \toprule
    \textbf{OpCode} & \textbf{Sequence} & \textbf{Comment} \\
    \midrule
    \endhead
    \midrule
    \multicolumn{2}{r}{\textit{Continued on next page}} \\
    \midrule
    \endfoot
    \bottomrule
    \endlastfoot
    
     % Combine the 6 branch instructions
    \textbf{BEQ} Rn, Rm, target & CBR.EQ Rn, Rm, target &
         \multirow{6}{*}{\parbox{1\linewidth}{subtract RegM from RegN and branch on condition.}} \\
    \textbf{BLT} Rn, Rm, target & CBR.LT Rn, Rm, target &  \\
    \textbf{BGT} Rn, Rm, target & CBR.GT Rn, Rm, target &  \\
    \textbf{BNE} Rn, Rm, target & CBR.NE Rn, Rm, target &  \\
    \textbf{BGE} Rn, Rm, target & CBR.GE Rn, Rm, target &  \\
    \textbf{BLE} Rn, Rm, target & CBR.LE Rn, Rm, target &  \\
    \midrule
    
     % Combine the 6 branch instructions
    \textbf{BEQ} Rn, target & MBR.EQ Rn, Rn, target &
         \multirow{6}{*}{\parbox{1\linewidth}{test RegN and branch on condition.}} \\
    \textbf{BLT} Rn, target & MBR.LT Rn, Rm, target &  \\
    \textbf{BGT} Rn, target & MBR.GT Rn, Rm, target &  \\
    \textbf{BEV} Rn, target & MBR.EV Rn, Rm, target &  \\
    \textbf{BNE} Rn, target & MBR.NE Rn, Rm, target &  \\
    \textbf{BGE} Rn, target & MBR.GE Rn, Rm, target &  \\
    \textbf{BLE} Rn, target & MBR.LE Rn, Rm, target &  \\
    \textbf{BOD} Rn, target & MBR.OD Rn, Rm, target &  \\
    \midrule
    
\end{longtable}





	